% Generated draft: requires mathematical and source review.
\documentclass[11pt]{article}
\usepackage[margin=25mm]{geometry}
\usepackage{fontspec}
\setmainfont{Latin Modern Roman}
\usepackage{amsmath,amssymb,mathtools,unicode-math}
\setmathfont{Latin Modern Math}
\usepackage{enumitem,needspace,xurl,hyperref}
\setlength{\parindent}{0pt}
\setlength{\parskip}{5pt}
\setlength{\emergencystretch}{3em}
\newcommand{\N}{\mathbb{N}}
\newcommand{\Z}{\mathbb{Z}}
\newcommand{\Q}{\mathbb{Q}}
\newcommand{\R}{\mathbb{R}}
\newcommand{\C}{\mathbb{C}}
\newcommand{\F}{\mathbb{F}}
\newcommand{\A}{\mathbb{A}}
\newcommand{\PP}{\mathbb{P}}
\newcommand{\E}{\mathbb{E}}
\newcommand{\eps}{\varepsilon}
\newcommand{\dd}{\,\mathrm{d}}
\begin{document}
\section*{UnsolvedMath ID: 800604}
\textbf{Source problem code:} catalogue entry
\subsection{Definitions and mathematical statement}
The supplied question concerns the lower bound
\[
\mathfrak{u}\leq \operatorname{add}\bigl(\mathsf{Split}(\mathrm{T},\mathrm{T})\bigr).
\]
This inequality is reported by the source record; it is not proved or independently verified here.

The ultrafilter number \(\mathfrak{u}\) is the least cardinality of a base for a free ultrafilter on \(\N\). Explicitly, writing \(B\subseteq^* A\) when \(B\setminus A\) is finite,
\[
\mathfrak{u}=\min\left\{|\mathcal{B}|:\begin{array}{l}
\mathcal{U}\text{ is a free ultrafilter on }\N,\quad \mathcal{B}\subseteq\mathcal{U},\\
\text{and }(\forall A\in\mathcal{U})(\exists B\in\mathcal{B})\ B\subseteq^* A
\end{array}\right\}.
\]
An ultrafilter is free if it contains no finite set.

The remaining notation requires conventions absent from the supplied record. Under a usual nontrivial-cover convention, a family \(\mathcal{V}\) of subsets of a space \(X\) is a \(\tau\)-cover if \(X\notin\mathcal{V}\), every point belongs to infinitely many members of \(\mathcal{V}\), and, for every \(x,y\in X\), at least one of
\[
\{V\in\mathcal{V}:x\in V,\ y\notin V\},
\qquad
\{V\in\mathcal{V}:y\in V,\ x\notin V\}
\]
is finite. The supplied record does not specify the admissible spaces or the openness, countability, or other restrictions defining its class \(\mathrm{T}(X)\); the preceding description therefore supplies a conventional combinatorial interpretation, not a verified identification of all the source's conventions.

For a specified class \(\mathrm{T}(X)\), the splitting property is expressed by
\[
X\models\mathsf{Split}(\mathrm{T},\mathrm{T})
\quad\Longleftrightarrow\quad
(\forall\mathcal{V}\in\mathrm{T}(X))
(\exists\mathcal{V}_0,\mathcal{V}_1\in\mathrm{T}(X))\
\bigl[\mathcal{V}=\mathcal{V}_0\mathbin{\dot\cup}\mathcal{V}_1\bigr].
\]
Once an ambient space \(Z\) and the cover conventions have been fixed, let
\[
\mathcal{P}_Z=\{X\subseteq Z:X\models\mathsf{Split}(\mathrm{T},\mathrm{T})\},
\]
with the subspace topology on each \(X\). The usual union additivity number is then
\[
\operatorname{add}(\mathcal{P}_Z)=
\min\left\{\kappa:\exists(X_\alpha)_{\alpha<\kappa}\in\mathcal{P}_Z^\kappa,
\ \bigcup_{\alpha<\kappa}X_\alpha\notin\mathcal{P}_Z\right\},
\]
provided a witnessing family exists. The record does not specify the intended ambient space or the convention when no witness exists.

\textbf{Source question (Problem 32).} Can the reported lower bound \(\mathfrak{u}\) on \(\operatorname{add}(\mathsf{Split}(\mathrm{T},\mathrm{T}))\) be improved?

The record names neither a proposed stronger bound nor the assumptions under which an improvement is sought. In particular, it does not explicitly ask for the strict inequality \(\mathfrak{u}<\operatorname{add}(\mathsf{Split}(\mathrm{T},\mathrm{T}))\). A fully specified conjecture requires clarification of these points and of the cover and ambient-space conventions.
\subsection{Short English statement}
Tsaban's Problem 32 asks whether the reported ultrafilter-number lower bound for the additivity of \(\mathsf{Split}(\mathrm{T},\mathrm{T})\) can be strengthened. The supplied record leaves the precise cover conventions, ambient class of spaces, and intended stronger bound unspecified. Its \emph{open} label is supplied registry status, not an independently verified assessment.
\subsection{Sources}
\begin{itemize}
\item Local UnsolvedMath registry, permanent ID \(800604\), source problem code \texttt{catalogue entry}; title: \emph{Selection Principles in Mathematics --- Problem 32 ([AddQuad])}. The record attributes the problem to Tsaban, with proposed year 2006, and labels it \emph{open}.
\item Tsaban, \emph{Selection Principles and Special Sets of Reals: Open Problems} (2006), Problem 32. Supplied source URL: \url{https://arxiv.org/abs/math/0609601}. The record reports access on 2026-10-16 and extraction by \texttt{arxiv\_latex}; these are supplied provenance claims, not evidence of an independent source visit here.
\item Supplied attribution/context label: \texttt{[AddQuad]}. No expansion or bibliographic identification of this label is included in the complete record.
\item Supplied upstream revision: \texttt{d371ee917ac9647fbb06cd6fc529deec6f3dda82}. The registry's update timestamp is \texttt{2026-10-16T14:00:00Z}; it is not an independently verified literature assessment.
\end{itemize}
\end{document}
